\scanpage{121}
% ms. line 01
\[ \sum n!\alpha_n(\lambda z)=\sum n!\alpha_n(\lambda)z^n \]
% ms. line 01a
converge
% ms. line 02
quel que soit $z$. Les deux \inserted{autres} relations
% ms. line 03
\deleted{\unsure{sont}} sont une autre façon d'exprimer la première,
% ms. line 04
-- en tenant compte de la formule de
% ms. line 05
Stirling $\sqrt[n]{n!}\sim\frac{n}{e}$ --.

% ms. line 06
Démonstration de la proposition. On peut
% ms. line 07
supposer les $\lambda_i>0$ \inserted{et \uline{décroissants}} $\unsure{]|;|.}$ soit alors $f(z)=\prod(1+\lambda_i z)$
% ms. line 08
\deleted{\unsure{une fonction}}, \inserted{$=\sum\alpha_n\deleted{(\lambda)}z^n$} \inserted{(où on pose $\alpha_n=\alpha_n(\lambda)$)} on a, \deleted{\unsure{en vertu de}} l'inégalité
% ms. line 09
de Cauchy
% ms. line 10
\[ \alpha_n\leq\frac{M(R)}{R^{n+1}}\qquad\text{pour tout }R>0 \]
% ms. line 11
où $M(R)$ est le maximum de $|f(z)|$ sur la
% ms. line 12
circonférence $|z|=R$, \unsure{d'où|d'ici|donc} $M(R)=\prod(1+\lambda_iR)$.
% ms. line 13
\deleted{Soit} On a donc
% ms. line 14
\[ \alpha_n\leq\left(\frac1R+\lambda_1\right)\ldots\left(\frac1R+\lambda_{n+1}\right)\prod_{i=n+2}^{\infty}(1+\lambda_i R) \]
% ms. line 15
\deleted{un facteur} \inserted{Or} $\prod_{i=n+2}^{\infty}(1+\lambda_i R)\leq e^{\rho_{n+1}R}$, où
% ms. line 16
$\rho_{n+1}=\sum_{k=n+2}^{\infty}\lambda_k$. faisons $R=\frac{n+1}{\rho_{n+1}}$, on obtient
% ms. line 17-18
\[ \begin{aligned} \alpha_n&\leq\left(\frac{\rho_{n+1}}{n+1}+\lambda_1\right)\ldots\left(\frac{\rho_{n+1}}{n+1}+\lambda_{n+1}\right)e^n \deleted{\unsure{\leq(\frac{\rho_1}{1}+\lambda_1)\ldots(\rho_{n+1}}}\\ &\leq\left(\frac{\unsure{\rho_1|\rho_2}}{1}+\lambda_1\right)\left(\frac{\rho_2}{2}+\lambda_2\right)\ldots\left(\frac{\rho_{n+1}}{n+1}+\lambda_{n+1}\right)e^n \end{aligned} \]
% ms. line 19
Or, \deleted{\unsure{comme les}} \inserted{les} $\lambda_i$ \deleted{\unsure{sont}} \unsure{étant|formant} décroissants et leur
% ms. line 20
suite sommable, on a $\lambda_n=\unsure{o|O}(\frac1n)$, d'autre part on
% ms. line 21
a \deleted{$\unsure{\rho_n\leq 0|\rho_n\to 0}$} $\frac{\rho_n}{n}=\uline{o}(\frac{1}{n})$, d'où $\frac{\rho_n}{n}+\lambda_n\leq\frac{\varepsilon_{n+1}}{n}$, où $\varepsilon_n\to0$
% ms. line 22
\[ \alpha_n\leq\frac{\varepsilon_{\unsure{1|0}}\varepsilon_2\ldots\varepsilon_{n+1}}{(n+1)!}e^n\qquad\text{où }\varepsilon_i\to0. \]
% ms. line 22a
D'où
% archive sheet number: 122
